\section{Introduction}\label{sec:intro}

Every driver who has squeezed into a tight space knows the trade-off: with one smooth reverse manoeuvre you need a generous gap, while a gap only slightly longer than your car can still be entered by shuffling back and forth many times.
This paper quantifies the trade-off.
For a kinematic car (front-wheel steering, fixed rear axle, bounded steering angle) we ask for the shortest gap
\[
  \Lmin(n)=\inf\{L:\ \text{the car can enter a gap of length $L$ using at most $n$ manoeuvres}\},
\]
where a \emph{manoeuvre} is a maximal stretch of motion in one direction (forward or reverse).

\paragraph{The model in brief.}
The car is a rectangle of length $\ell=d_f+d_w+d_b$ and width $w$.
Its reference point is the midpoint of the rear axle, and that point moves along curves of curvature at most $1/R$, where $R=d_w\cot\theta$ is the minimal turning radius.
The curb is the line $y=0$ and is a rigid obstacle for the whole car.
The two parked cars have the same width $w$ and stand at distance $m$ from the curb, so they occupy the strip $m\le y\le H:=m+w$.
One ends at $x=0$ and the other begins at $x=L$.
Our car has \emph{parked} when it is collision-free and lies entirely in $0\le y\le H$, i.e.\ its outermost point is no farther out than the outer sides of the parked cars.
Throughout, the parameters are fixed and the main application is $m=w/10$.
Precise definitions are given in \cref{sec:model}, and \cref{fig:geometry} shows the geometry.

\paragraph{Notation used below.}
$b=d_w+d_f$ is the distance from the rear axle to the front bumper and $\Rpm=R\pm w/2$.
Further, $\rho=\sqrt{b^2+\Rp^2}$ and $r=\sqrt{d_b^2+\Rp^2}$ are the distances from the left turning centre to the two curb-side corners.

\subsection{Main results}

\paragraph{One manoeuvre.}
Driving out of a parked position is the time reversal of parking, so we may study \emph{exits}.
In the optimal single manoeuvre the car starts from the parked position that is as far back and as far out as possible, tilted nose-up by a small angle $\psi_d$.
It then drives forward with full left lock, and its front curb-side corner grazes the outer corner of the car in front.
The tilt is limited by the curb: the rear overhang swings towards the curb during the turn and may just touch it.

\begin{theorem}[one manoeuvre; precise version \cref{thm:one}]\label{thm:intro-one}
Assume the mild geometric conditions of \cref{ass:one}, which hold for typical passenger cars when $m=w/10$ (\cref{app:assumptions}).
Let $\varsigma=r-H$, $s_0^2=b^2+\Rm^2$, and let $\psi_d\in[0,\pi/2)$ solve $\Rm\cos\psi_d-b\sin\psi_d=\varsigma$.
Then
\begin{equation}\label{eq:Lam1}
\begin{aligned}
  \Lmin(1)=\Lambda_1&:=d_b\cos\psi_d-\Rm\sin\psi_d+\sqrt{\rho^2-\varsigma^2}\\
  &=\frac{\ell\Rm \varsigma+(bd_b-\Rm^2)\sqrt{s_0^2-\varsigma^2}}{s_0^2}+\sqrt{\rho^2-\varsigma^2}.
\end{aligned}
\end{equation}
\end{theorem}

When the margin equals the swing $\varepsilon_b:=r-\Rp$ of the rear overhang, $\psi_d=0$ and \eqref{eq:Lam1} becomes the classical one-manoeuvre length $d_b+\sqrt{b^2+2Rw}$; the term $\sqrt{b^2+2Rw}$ is the formula usually attributed to Blackburn \cite{Blackburn}, which measures the gap from the rear axle.
A larger margin allows a tilted final position and a slightly shorter slot.

\paragraph{Many manoeuvres.}

\begin{theorem}[\cref{thm:mono,thm:asym}]\label{thm:intro-asym}
$\Lmin(n)$ is non-increasing in $n$. Under the hypotheses below it satisfies $\ell<\Lmin(n)\le\Lmin(1)$ for all $n$, and $\Lmin(n)\to\ell$.
More precisely, for every car with $R>w/2$, $d_b\le b$, $\ell\ge2w$ and every margin $0<m\le w/2$ (a convenient special case of the hypotheses of \cref{thm:asym}),
\[
  \lim_{n\to\infty}\sqrt n\,\big(\Lmin(n)-\ell\big)=K,\qquad K=\sqrt{4Rw-\frac{w^3}{6R}}.
\]
Equivalently, a slot with slack $g=L-\ell\to0$ requires $(K^2+o(1))/g^2$ manoeuvres, and this many suffice.
\end{theorem}

The constant is remarkably universal: it involves only the turning radius and the width.
The asymptotically optimal strategy is a sideways \emph{crawl}.
The car bounces between the two parked cars, and each manoeuvre is an S-curve made of two full-lock arcs that lifts the car by about $g^2/(4R)$ away from the curb.
Once the car is more than halfway out, the S-curves are slightly tilted.
That tilt lets the overhangs pass over the parked cars and is responsible for the correction $-w^3/(6R)$.
The matching lower bound comes from a confinement lemma together with a \emph{calibration}, a potential function whose increment bounds the lateral progress of any single manoeuvre (\cref{sec:asym}).

\paragraph{Few manoeuvres.}
Several manoeuvres pay off already for small $n$.
As an exit, the best simple two-manoeuvre strategy first reverses with full right lock until the rear corners touch the curb and the car behind simultaneously, and then drives out with full left lock.
This gives the closed-form bound $\Lmin(2)\le\Lambda_2$ of \cref{prop:two}.
Starting with the front bumper against the car in front and first reversing briefly with left lock does slightly better (\cref{prop:two-better}).
\Cref{tab:values} lists our values for a typical car.

\begin{table}[ht]
\centering
\begin{tabular}{lccccc}
\toprule
$n$ & $1$ & $2$ & $3$ & $4$ & $\infty$\\
\midrule
$\Lmin(n)$ [m], rigid curb & $5.833$ & $\le5.411$ & $\le5.280$ & $\le5.180$ & $\ell=4.300$\\
$\Lmins(n)$ [m], soft curb & $5.809$ & $\le5.312$ & $\le5.088$ & $\le4.978$ & $\ell=4.300$\\
\bottomrule
\end{tabular}
\caption{Typical car: $d_f=0.9$, $d_w=2.6$, $d_b=0.8$, $w=1.8$ (metres), $\theta=35^\circ$, margin $m=0.18$.
Here $R=3.713$, the classical one-manoeuvre value is $5.861$, and $K=5.145$.
For $n=1$ the values are exact, and for $n\ge2$ they come from explicit strategies (\cref{sec:several,sec:soft}) and are rounded up. Other car sizes are in \cref{tab:cars,tab:compare}.}
\label{tab:values}
\end{table}

\paragraph{A soft curb.}
Ordinary curbs are low enough for the bumpers to pass over them, and then only the tyres must stay on the road.
We treat this \emph{soft-curb} variant in \cref{sec:soft} and compare the two models in \cref{sec:compare}.
The optimal single manoeuvre keeps its shape and its closed form, with a larger admissible tilt; for the typical car it needs $5.809$ m.
For two manoeuvres the textbook exit becomes admissible: reverse with full right lock until the rear bumper touches the car behind, then leave with full left lock.
This exit needs $5.312$ m (\cref{prop:two-soft}).
For two to five manoeuvres our soft-curb strategies are $10$--$20$ cm shorter than the rigid-curb ones for the typical car, and up to about $40$ cm shorter for larger cars, while the asymptotic constant $K$ is the same (\cref{prop:asym-soft}).

\paragraph{Why a margin?}
The question as usually posed requires the parked car to lie in a slot exactly as deep as the car is wide, i.e.\ flush with the curb.
With a rigid curb that problem is degenerate: a car with positive rear overhang can never become flush with a straight wall by a finite motion, whatever the slot length and the number of manoeuvres (\cref{prop:flush}).
The margin $m>0$ removes this degeneracy, and it is also what real parking requires.

\subsection{Related work}
The one-manoeuvre problem was popularised by Blackburn \cite{Blackburn}; up to the rear overhang, his formula is the case $m=\varepsilon_b$ of \cref{thm:intro-one}.
Engineering studies of automatic parallel parking compute minimal slot dimensions for specific families of circular-arc paths with several trials \cite{WangEtAl,Vorobieva}.
They observe that the required space tends to the size of the vehicle as the number of trials grows, but they do not prove optimality or rates.
The kinematic car is the Dubins car \cite{Dubins} when it only moves forward, and the Reeds--Shepp car \cite{ReedsShepp} when reversals are allowed; the latter is a classical example of a nonholonomic system that is small-time locally controllable.
Laumond's theorem \cite{Laumond} shows that it can follow any collision-free path approximately, which explains qualitatively why $\Lmin(n)\to\ell$.
The contribution here is quantitative and sharp: exact values for one manoeuvre, and the exact first-order asymptotics $K/\sqrt n$.

\subsection{Outline}
\Cref{sec:model} defines the model and proves elementary properties.
\Cref{sec:toolkit} collects the geometric tools: turning centres, the drift of the centres, and a distance identity.
\Cref{sec:one} proves \cref{thm:one} and discusses the zero-margin degeneracy.
\Cref{sec:several} treats few manoeuvres, and \cref{sec:asym} proves the asymptotic theorem.
\Cref{sec:soft} treats the soft curb, and \cref{sec:compare} compares the two models.
\Cref{sec:lean} describes the formal verification in Lean.
\Cref{sec:discussion} discusses variants and open problems.
The appendix records the verification of the geometric assumptions, values for other cars, five- and twenty-manoeuvre parking motions with a rigid and with a soft curb, our numerical methods, and a disclosure of how an LLM was used.
